\clearpage
\VerdierPhysicalPageAnchor{0247}%
\begingroup\small\setlength{\parskip}{0pt}\setlength{\emergencystretch}{1.5em}
\setlength{\abovedisplayskip}{4pt plus 2pt minus 2pt}%
\setlength{\belowdisplayskip}{4pt plus 2pt minus 2pt}%

\VerdierSemanticAnchor{ch4d-4-6-5-resolution-type}
\noindent 是内射型分解 (1.1.4)。

\VerdierSemanticAnchor{ch4d-4-6-5-terminology}
因此，我们所谓的\emph{复形 $Y$ 的 Cartan--Eilenberg 内射分解}，就是这些作者
简称为\emph{复形 $Y$ 的内射分解}的对象（[1]，第 XVII 章）；本书不采用
后一术语。$Y$ 的一个 Cartan--Eilenberg 内射分解诱导出复形 $Y$ 各分量
以及 $Y$ 各微分之核的内射型分解［\emph{同上引文}］。

\VerdierSemanticAnchor{ch4d-4-6-5-finite-length}
若存在整数 $n_0$，使得当 $n>n_0$ 时 $I^{\bullet,n}=0$，则称
Cartan--Eilenberg 分解
$\varepsilon\colon Y\longrightarrow I^{\bullet\bullet}$ 具有
\emph{有限长度}。复形 $Y$ 存在有限长度的 Cartan--Eilenberg 内射分解的
充要条件是：$Y$ 的上同调对象和 $Y$ 的各微分之像的内射维数有上界。
回顾\emph{同上引文}中建立的结果：

\VerdierSemanticAnchor{ch4d-4-6-6-a}
\medskip\noindent\textbf{命题 4.6.6.}\quad
\textit{a) 当范畴 $\mathcal A$ 有足够多的内射对象 (3.1.5) 时，
$\mathcal A$ 的每个单复形 $Y$ 都有 Cartan--Eilenberg 内射分解。}

\VerdierSemanticAnchor{ch4d-4-6-6-b}
\noindent\textit{b) 设
$f\colon X\longrightarrow Y$ 是单复形的一个态射，
$\varepsilon\colon X\longrightarrow I^{\bullet\bullet}$ 和
$\varepsilon'\colon Y\longrightarrow J^{\bullet\bullet}$ 是两个
Cartan--Eilenberg 内射分解。存在双复形的一个态射}

\VerdierSemanticAnchor{ch4d-4-6-6-b-morphism}
\[
\textit{$g\colon I^{\bullet\bullet}\longrightarrow J^{\bullet\bullet}$}
\]
\textit{位于 $f$ 之上；即双复形态射的图}

\VerdierSemanticAnchor{ch4d-4-6-6-b-diagram}
\[
\begin{tikzcd}[column sep=5.0em,row sep=3.0em]
X \arrow[r,"f"] \arrow[d,"\varepsilon"'] &
Y \arrow[d,"\varepsilon'"] \\
I^{\bullet\bullet} \arrow[r,"g"'] &
J^{\bullet\bullet}
\end{tikzcd}
\]
\textit{交换。}

\VerdierSemanticAnchor{ch4d-4-6-6-c}
\noindent\textit{c) 位于 $f$ 之上的两个态射
$g,g'\colon I^{\bullet\bullet}\longrightarrow J^{\bullet\bullet}$ 同伦
（这里所用的是双复形态射的同伦，见第一章，2.5.1）。}

\VerdierSemanticAnchor{ch4d-4-6-6-proof}
证明参见\emph{同上引文}。

\VerdierSemanticAnchor{ch4d-4-6-7}
\medskip\noindent\textbf{4.6.7.}\quad
设 $Y$ 是 $\operatorname{comp}(\mathcal A)$ 的一个对象，且
$Y\longrightarrow I^{\bullet\bullet}$ 是 $Y$ 的一个
Cartan--Eilenberg 内射分解。双复形 $I^{\bullet\bullet}$ 按其第一次数或
第二次数滤过时，分别产生两个取值于 $D(\mathcal A)$ 的谱对象 (4.1.7)。
由按第二次数滤过的 $I^{\bullet\bullet}$ 所定义的、取值于
$D(\mathcal A)$ 的谱对象称为\emph{$Y$ 的 Cartan--Eilenberg 谱对象}
（相对于所考察的 Cartan--Eilenberg 分解）。
\endgroup
